% A+B paper: the transfer study (spine) + matched-baseline negative result & Test D mechanism (core).
% Scaffold generated 2026-07-20. Reuses references.bib. Related Work / Method prose can be ported
% from paper.tex (archived old paper). Numbers are from notes/EXPERIMENT_LOG.md.
% DO NOT write toward the compositional-prior claim (C) until the Spider+execution result lands.
\documentclass[11pt]{article}
\usepackage[utf8]{inputenc}
\usepackage{amsmath,amssymb}
\usepackage{booktabs}
\usepackage{graphicx}
\usepackage{hyperref}
\usepackage{natbib}
\usepackage{geometry}
\usepackage{enumitem}
\usepackage{xcolor}
\usepackage{listings}

\geometry{margin=1in}
\newcommand{\todo}[1]{\textcolor{red}{[TODO: #1]}}

\title{What Transfers from Text-to-SQL to Text-to-Cypher?\\A Controlled Study of Reasoning, Selection, and Training Recipe}
\author{Cyrus Parsons}
\date{}

\begin{document}
\maketitle

\begin{abstract}
Text-to-SQL has produced a large toolbox of techniques---chain-of-thought (CoT) distillation,
self-consistency, execution-based selection, reinforcement learning---but it is unclear which of
them transfer to the structurally different task of Text-to-Cypher. We run a controlled,
matched-pipeline study to find out, unified by a single \emph{constrained-output-space hypothesis}:
a Cypher graph pattern is a single \emph{connected} object with few semantically-equivalent surface
forms, whereas SQL is \emph{compositional} with many. This predicts that (P1)~string-diversity
methods fail, (P2)~execution-grounded methods still transfer, and (P3)~CoT/decomposition hurts.
Our central result is negative and, to our knowledge, the first controlled demonstration of it: CoT
distillation---the highest-profile SQL transfer---does \emph{not} improve Text-to-Cypher and
modestly hurts it. The effect holds across two model families (Gemma-2-9B, Llama-3.1-8B), naive and
execution-verified traces, and in-distribution, compositional, and length-generalization regimes.
Apparent gains reported under an earlier framing are attributable to a stronger supervised
fine-tuning recipe and a leaked benchmark split, not to reasoning. We give a causal mechanism: CoT's
decompositional bias fragments connected graph patterns, and a holistic-reasoning ablation recovers
$\sim$40\% of the penalty by cutting fragmentation 5--6$\times$. Consistent with the hypothesis,
execution-grounded selection transfers (+1.1pp exact match) while string self-consistency does not.
We close with a methodological lesson: a published artifact used as a control, rather than a matched
in-pipeline baseline, produced months of false ``CoT helps'' conclusions.
\end{abstract}

%------------------------------------------------------------------
\section{Introduction}
%------------------------------------------------------------------
Text-to-SQL is a mature field with a well-stocked toolbox: CoT distillation \citep{he2025starsql,
tai2023cot}, self-consistency, execution-based candidate selection, and reinforcement learning have
each produced large gains. Text-to-Cypher---translating natural language to the graph query language
of Neo4j \citep{ozsoy2025text2cypher}---is younger and structurally different, and it is tempting to
assume the SQL toolbox carries over. This paper asks, in a controlled setting, \emph{which techniques
actually transfer, and why}.

Our organizing idea is the \textbf{constrained-output-space hypothesis} (\S\ref{sec:hyp}): a Cypher
basic graph pattern is a single \emph{connected} object admitting few semantically-equivalent surface
forms, unlike SQL, which composes from sub-queries and joins and admits many. This structural
difference makes three predictions about transfer (P1--P3 above).

Our headline finding tests the most important transfer---CoT distillation---and it is negative.
Under a matched pipeline where the \emph{only} variable is the training target (a reasoning trace
followed by the query, versus the query alone), CoT distillation modestly \emph{degrades}
Text-to-Cypher across two model families and every regime we test. We trace an earlier positive
framing of this same project to two confounds: an unmatched baseline (a published adapter used as a
control) and a leaked cross-split evaluation. We then explain the negative result mechanistically and
show which techniques \emph{do} transfer.

\noindent\textbf{Contributions.}
\begin{enumerate}[leftmargin=*]
  \item \textbf{A controlled matched-baseline result} that CoT distillation does not help, and
  modestly hurts, Text-to-Cypher---across Gemma-2-9B and Llama-3.1-8B, naive and execution-verified
  traces, and in-distribution / compositional / length-generalization regimes (\S\ref{sec:matched}).
  \item \textbf{A causal mechanism}: decomposition fragments connected graph patterns; a
  holistic-reasoning ablation recovers $\sim$40\% of the penalty (\S\ref{sec:mech}).
  \item \textbf{A positive transfer}: execution-grounded (MBR) selection beats string voting while
  string self-consistency fails, confirming the hypothesis (\S\ref{sec:positive}).
  \item \textbf{A methodological lesson}: always train a matched in-pipeline baseline before
  attributing a delta to an intervention (\S\ref{sec:discussion}).
\end{enumerate}

%------------------------------------------------------------------
\section{Related Work}
%------------------------------------------------------------------
\todo{Port from paper.tex \S Related Work (Text2Cypher; CoT-for-SQL incl.\ STaR-SQL \citep{he2025starsql}
and \citet{tai2023cot}; RL \citep{tran2025refining}; selection; benchmarks). Reframe: these are the
SQL toolbox whose transfer we test. Emphasize that CoT-helps-SQL is an \emph{assumed-positive} we
overturn in a controlled setting for Cypher.}

%------------------------------------------------------------------
\section{The Constrained-Output-Space Hypothesis}
\label{sec:hyp}
%------------------------------------------------------------------
\todo{Formalize. A Cypher pattern \texttt{(a)-[:R]->(b)-[:S]->(c)} is one connected object: its
sub-parts cannot be assembled independently without a global consistency constraint (shared
variables). SQL composes from sub-queries/CTEs/joins that \emph{can} be assembled independently, and
admits many equivalent surface forms (JOIN vs subquery, alias renaming, clause order). Derive
predictions:}
\begin{itemize}[leftmargin=*]
  \item \textbf{P1 (string diversity fails).} With ~one canonical form, sampling produces little
  useful diversity, so string-vote self-consistency cannot help.
  \item \textbf{P2 (execution grounding transfers).} Methods that operate on \emph{results} bypass
  the surface-form bottleneck and still help.
  \item \textbf{P3 (decomposition hurts).} Decomposing a connected pattern fragments it; the right
  inductive bias for SQL is the wrong one for Cypher.
\end{itemize}

%------------------------------------------------------------------
\section{Matched Experimental Protocol}
\label{sec:protocol}
%------------------------------------------------------------------
Every comparison in this paper fixes the base model, the QLoRA configuration (4-bit NF4, $r{=}64$,
$\alpha{=}64$, all-linear targets, \texttt{paged\_adamw\_8bit}), the completion-only loss masking,
and the CoT-distillation procedure (traces from GPT-oss-120B), and varies \emph{only} the training
target: the reference query alone (\emph{direct}) versus a reasoning trace followed by the reference
query (\emph{CoT}). Both arms train on the identical set of instances (those with a valid trace), so
no data differs between them. We instantiate this protocol on Cypher (Neo4j 2024; ZOGRASCOPE),
SPARQL (LC-QuAD 2.0), and SQL (\S\ref{sec:crossform}).
\todo{Port dataset/QLoRA/inference details from paper.tex \S Method + \S Experimental Setup.}

%------------------------------------------------------------------
\section{Results}
%------------------------------------------------------------------

\subsection{The matched negative result on Text-to-Cypher}
\label{sec:matched}
Table~\ref{tab:matched} reports the matched comparison on the Neo4j 2024 benchmark. For \emph{both}
model families, CoT \emph{lowers} every translation and execution metric.

\begin{table}[t]\centering
\caption{Matched direct-vs-CoT on Neo4j 2024 (our pipeline; only the training target differs).
CoT hurts both families.}
\label{tab:matched}
\begin{tabular}{llccc}
\toprule
Model & Target & GLEU & String EM & Exec EM \\
\midrule
Gemma-2-9B & direct & \textbf{0.7854} & \textbf{0.4331} & --$^{\dagger}$ \\
Gemma-2-9B & CoT    & 0.7682 & 0.3799 & 0.2554 \\
\addlinespace
Llama-3.1-8B & direct & \textbf{0.7680} & \textbf{0.4223} & \textbf{0.2865} \\
Llama-3.1-8B & CoT    & 0.7416 & 0.3000 & 0.2161 \\
\midrule
\multicolumn{2}{l}{CoT effect (Gemma)} & $-0.0172$ & $-0.0532$ & -- \\
\multicolumn{2}{l}{CoT effect (Llama)} & $-0.0264$ & $-0.1223$ & $-0.0704$ \\
\bottomrule
\end{tabular}
\\[2pt]{\footnotesize $^{\dagger}$Gemma direct-answer execution EM pending (predictions generated).}
\end{table}

The positive result reported under an earlier framing of this project decomposes cleanly
(Table~\ref{tab:decomp}): the apparent ``+0.1227 GLEU from CoT'' was a comparison against a
\emph{published} adapter run through our inference harness, not a matched control. Replacing that
baseline with our own direct-answer model shows the gain was the training pipeline ($+0.1399$); CoT
itself \emph{subtracts} ($-0.0172$).

\begin{table}[t]\centering
\caption{Decomposition of the previously reported ``+0.1227 GLEU from CoT'' on Neo4j 2024.}
\label{tab:decomp}
\begin{tabular}{lc}
\toprule
Component & $\Delta$ GLEU \\
\midrule
Pipeline (Neo4j published adapter $\to$ our direct-answer adapter, same inference) & $+0.1399$ \\
CoT (our direct-answer $\to$ our CoT) & $-0.0172$ \\
\midrule
Net (previously reported as ``CoT'') & $+0.1227$ \\
\bottomrule
\end{tabular}
\end{table}

\subsection{Robustness: leakage, verified traces, and generalization splits}
\label{sec:robust}
\textbf{Not a leakage artifact.} Neo4j 2024 contains 31.7\% of test questions verbatim in training,
but the leakage is adversarial (often a \emph{different} gold query); on genuinely unseen questions
the direct-answer model beats CoT by $+0.068$ EM, so CoT loses \emph{most} where leakage is absent.
\textbf{Verified traces don't rescue it.} Retraining both arms on execution-verified forward traces
(the STaR-SQL ingredient) still yields CoT below direct (exec EM $0.0939$ vs $0.1052$).
\textbf{Generalization.} On clean, per-split ZOGRASCOPE (leakage-verified), direct-answer wins every
split including length (Table~\ref{tab:zog})---overturning an earlier ``length-generalization SOTA''
claim that was a cross-split leakage artifact.

\begin{table}[t]\centering
\caption{Clean ZOGRASCOPE execution accuracy (per-split, leakage-verified). Direct wins every split.}
\label{tab:zog}
\begin{tabular}{lccc}
\toprule
Split & Direct & CoT & $\Delta$ (direct $-$ CoT) \\
\midrule
IID            & 0.8255 & 0.6094 & $+0.216$ \\
Compositional  & 0.6612 & 0.4181 & $+0.243$ \\
Length (dist.\ shift) & 0.4397 & 0.2306 & $+0.209$ \\
\bottomrule
\end{tabular}
\end{table}

\subsection{Cross-formalism}
\label{sec:crossform}
Table~\ref{tab:crossform} extends the matched comparison to SPARQL and SQL. CoT hurts SPARQL (GLEU)
as it does Cypher. On synthetic SQL, CoT also hurts under string and SQL-canonicalized exact match;
however, string metrics undercount SQL correctness (many equivalent forms), so the decisive test is
\emph{execution accuracy on Spider}, which is pending.
\todo{Fill the Spider execution row when the run lands. Per Alex, do not draw the compositional-prior
conclusion (C) from this table until that number exists.}

\begin{table}[t]\centering
\caption{Matched CoT effect across formalisms (our pipeline). Negative everywhere measured; the
decisive SQL execution-accuracy control is pending.}
\label{tab:crossform}
\begin{tabular}{lll}
\toprule
Formalism (dataset) & Metric & CoT effect \\
\midrule
Cypher (Neo4j 2024) & GLEU & $-0.0172$ \\
SPARQL (LC-QuAD 2.0) & GLEU & $-0.1470$ \\
SQL (synthetic) & canon.\ EM & $-0.0477$ \\
SQL (Spider) & execution acc. & \emph{pending} \\
\bottomrule
\end{tabular}
\end{table}

\subsection{What does transfer: execution-grounded selection and a stronger recipe}
\label{sec:positive}
Consistent with P1--P2, string-vote self-consistency falls \emph{below} greedy, while
execution-result (MBR) selection rises above both (Table~\ref{tab:select}); the large oracle ceiling
motivates a trained verifier as future work. Separately, completion-only loss masking yields a
direct-answer model well above the published baseline (GLEU 0.7854 vs 0.5560), of which the masking
choice accounts for $+0.044$ (a full-sequence ablation).

\begin{table}[t]\centering
\caption{Self-consistency selection on the CoT model (Neo4j exec subset). Execution-result voting
transfers; string voting does not.}
\label{tab:select}
\begin{tabular}{lc}
\toprule
Selection method & Exec EM \\
\midrule
String-vote SC@5 & 0.2509 \\
Greedy (single) & 0.2554 \\
Execution-vote SC@5 (MBR) & \textbf{0.2665} \\
Oracle (any of 5 correct) & 0.4302 \\
\bottomrule
\end{tabular}
\end{table}

%------------------------------------------------------------------
\section{Mechanism: why CoT hurts Cypher}
\label{sec:mech}
%------------------------------------------------------------------
On clean ZOGRASCOPE with execution ground truth, CoT's failures are compositional, not relational:
97\% of failures use the correct relationships, but assemble them into the wrong connected structure.
A reasoning-format ablation ladder (Table~\ref{tab:ladder}) isolates the cause. A \emph{holistic-path}
format that states the full connected path without sub-question decomposition recovers $\sim$40\% of
the penalty and cuts path fragmentation 5--6$\times$; adding explicit hop-enumeration (E4) does not
help further. The residual gap to direct-answer reflects a planning-stage truncation on long paths
(the reasoning itself commits to a too-short path 99--100\% of the time) plus an intrinsic cost of
generating reasoning before a jointly-constrained object.

\begin{table}[t]\centering
\caption{Reasoning-format ablation (ZOGRASCOPE, execution accuracy, common instances). Holistic
framing recovers part of the CoT penalty; no format closes the gap to direct.}
\label{tab:ladder}
\begin{tabular}{lccc}
\toprule
Config & Length & IID & Comp. \\
\midrule
Direct-answer & \textbf{0.475} & \textbf{0.826} & \textbf{0.661} \\
QDecomp-CoT   & 0.258 & 0.609 & 0.418 \\
Holistic-CoT  & 0.303 & 0.710 & 0.511 \\
Enum-CoT (E4) & 0.231 & 0.669 & 0.429 \\
\bottomrule
\end{tabular}
\end{table}

%------------------------------------------------------------------
\section{Discussion}
\label{sec:discussion}
%------------------------------------------------------------------
\todo{(1) The constrained-output-space account ties P1--P3 together: the same property that makes
string diversity useless (P1) and CoT harmful (P3) is what makes execution grounding the only
transfer that survives (P2). (2) Methodological lesson: a published artifact is not a control; the
matched baseline reversed the sign of our headline. (3) Scope: why SQL's CoT gains need not transfer;
flag the Spider+execution control as the test that would license a compositional-prior claim
(future work, not a claim here).}

%------------------------------------------------------------------
\section{Limitations}
%------------------------------------------------------------------
\todo{String metrics undercount SQL (motivates the pending Spider execution control); single teacher
(GPT-oss-120B); 2024-dataset leakage (which we turn into evidence, \S\ref{sec:robust}); approximate
result-set comparison for execution accuracy.}

%------------------------------------------------------------------
\section{Conclusion}
%------------------------------------------------------------------
\todo{One paragraph: which SQL techniques transfer to Cypher (execution-grounded selection; a strong
SFT recipe) and which do not (CoT distillation; string diversity), unified by the
constrained-output-space hypothesis and explained mechanistically.}

\bibliographystyle{plainnat}
\bibliography{references}
\end{document}
